\documentclass[10pt,a4paper]{amsart}
\usepackage[margin=19mm]{geometry}
\usepackage{amsmath,amssymb,mathtools}
\usepackage[scr=rsfs,cal=euler]{mathalfa}
\usepackage{xfrac}
\usepackage[unicode,hidelinks,bookmarks=false]{hyperref}
\newcommand{\ie}{i.e.\ }
\newcommand{\cf}{cf.\ }
\newcommand{\resp}{respectively\ }
\newcommand{\Subsection}{\subsection}
\providecommand{\nobreakdash}{\nobreak\hbox{-}}
\setlength{\emergencystretch}{2em}
\multlinegap0pt
\usepackage{bm}
\usepackage{bbm}
\usepackage{dsfont}
\usepackage{xspace}
\usepackage{cancel}
\usepackage{mathtools}
\usepackage{amsthm}
\makeatletter
\@ifundefined{lemma}{\newtheorem{lemma}{Lemma}}{}
\@ifundefined{theorem}{\newtheorem{theorem}[lemma]{Theorem}}{}
\makeatother
\title[Weight-constrained cut-down de Bruijn sequences]{Weight-constrained cut-down de Bruijn sequences}
\author{Chris J Mitchell, Peter R Wild}
\date{Source: arXiv:2607.22350v2 (CC BY 4.0)}
\begin{document}
\maketitle

\noindent\emph{Source and licence.} Weight-constrained cut-down de Bruijn sequences. Version arXiv:2607.22350v2 (CC BY 4.0), distributed under the Creative Commons Attribution 4.0 International licence, as recorded in the arXiv licence metadata for this exact version and shown on the abstract page https://arxiv.org/abs/2607.22350v2. Full rights notice: licenses/arxiv-2607.22350-CC-BY-4.0.txt.\par\smallskip

\noindent\textit{Editorial scope.} Counted unit: the Theorem whose source label is \texttt{theorem:P\_characteristic\_a} in the article (source0.tex lines 731--739, source label \texttt{theorem:P\_characteristic\_a}), delivered together with the article's own complete original proof of it (source0.tex lines 741--795, one \texttt{proof} environment of 55 lines). No mathematics is written, repaired or re-derived: statement and proof are the article's own text line for line. Required context retained uncounted: lemma:P\_balanced (lines 696--699). Every definition, display and label of those lines is retained unchanged. Omitted from this package and not redistributed: 1--695, 700--730, 740--740, 796--1029. Nothing inside a retained excerpt is omitted or altered: every retained body is delivered exactly as the article's lines are written, so no omission receipt of any kind accompanies this package. Citations: the retained excerpts carry no citation command at all, so no bibliography entry of the article is transcribed and no reference is redistributed. Reference adaptation: none was needed and none was made; every reference inside the delivered unit resolves inside it, so no unresolved reference or citation is produced. Printed slips kept literal: no printed word, symbol, hypothesis or number of the counted statement or of its proof was altered. Layout and class: the article's own class is \texttt{article} with packages including graphicx, amssymb, amsthm, subcaption, amsfonts, amsmath, hyperref, xcolor; the wrapper is the lane's pinned \texttt{amsart} editorial wrapper at 10pt on A4, which restates the theorem-like environments the retained text needs and adds the article's own macro definitions that the retained text needs (none), with extra wrapper packages (bm, bbm, dsfont, xspace, cancel, mathtools). The delivered numbering is the unit's own. Assets: no retained span contains a figure environment, a graphics-inclusion command, a PDF-page inclusion, a table or a TikZ picture; no image file is copied into this package, the package declares no asset dependency and no image file is redistributed with it. Licence: version arXiv:2607.22350v2 of this article is distributed under the Creative Commons Attribution 4.0 International licence, as recorded in the arXiv licence metadata for this exact version and shown on the abstract page https://arxiv.org/abs/2607.22350v2. A full rights notice is delivered at licenses/arxiv-2607.22350-CC-BY-4.0.txt. Characters: every retained excerpt is pure ASCII, so the delivered engine, XeLaTeX with its default Latin Modern text font, reports no missing character.
\begin{lemma}  \label{lemma:P_balanced}
If $k>2$, $n\geq 2$ and $n\leq r\leq s\leq n(k-1)$ (where $r,s\in\mathbb{Z}$ if $k$ is even and
$r=t/2$ and $s=u/2$ for some $t,u\in\mathbb{Z}$ if $k$ is odd) then $P_k^{r,s}(n-1)$ is balanced.
\end{lemma}

\begin{theorem}  \label{theorem:P_characteristic_a}
Suppose $n\geq2$, $k>2$, and $n\leq r\leq n(k-1)$ (where $r\in\mathbb{Z}$ if $k$ is even and
$r=t/2$ for some $t\in\mathbb{Z}$ if $k$ is odd). Then $P_k^{r,r}(n-1)$ is Eulerian if and only if
\begin{itemize}
\item[(i)] $r\in\{n,(k-1)n\}$,
\item[(ii)] $r\in\{n+1,(k-1)n-1\}$ and $k\neq3$, or
\item[(iii)]  $r\in\{n-1+k/2,(k-1)(n-1)+k/2\}$ and $k$ odd.
\end{itemize}
\end{theorem}

\begin{proof}
Because of Lemma~\ref{lemma:P_balanced}, to show the Eulerian property holds it suffices to show
that the relevant subgraph is strongly connected.
\begin{itemize}
\item[(i)] $P_k^{n,n}(n-1)$ consists of a single edge, namely $1^n$ --- a loop from the vertex
    $1^{n-1}$ to itself; this is clearly Eulerian.  Similarly, the subgraph
    $P_k^{(k-1)n,(k-1)n}(n-1)$ consists of the single edge $(k-1)^n$ and thus is also Eulerian.
\item[(ii)] $P_4^{n+1,n+1}(n-1)$ contains all the $n$-tuples in the two necklaces $[1^{n-1},2]$
    and $[1^{n-1},0]$.  These two necklaces share a vertex, namely $1^{n-1}$, and hence the
    subgraph is strongly connected. If $k>4$, then $P_k^{n+1,n+1}(n-1)$ contains only a single
    necklace, namely $[1^{n-1},2]$, and thus the subgraph is strongly connected.  Analogous
    arguments hold for $P_4^{3n-1,3n-1}(n-1)$, which contains the necklaces $[3^{n-1},2]$ and
    $[3^{n-1},0]$, and $P_k^{n(k-1)-1,n(k-1)-1}(n-1)$, which contains the single necklace
    $[(k-1)^{n-1},k-2]$.
\item[(iii)] If $k$ is odd then $n-1+k/2$ is not an integer, and hence the only necklace in
    $P_k^{n-1+k/2,n-1+k/2}(n-1)$ is $[1^{n-1},0]$, and so $P_k^{n-1+k/2,n-1+k/2}(n-1)$ is
    Eulerian.  Similarly, $(k-1)(n-1)+k/2$ is not an integer, and hence the only necklace in
    $P_k^{(k-1)(n-1)+k/2,(k-1)(n-1)+k/2}(n-1)$ is $[(k-1)^{n-1},0]$ and so
    $P_k^{(k-1)(n-1)+k/2,(k-1)(n-1)+k/2}(n-1)$ is Eulerian.
\end{itemize}
It remains to show that in all other cases $P_k^{r,r}(n-1)$ is not Eulerian. First consider the
special case $k=3$ and $r\in\{n+1,(k-1)n-1\}$. Observe that $P_3^{n+1,n+1}$ contains two necklaces:
$[1^{n-1},2]$ and $[1^{n-2},0,0]$. These clearly do not share a vertex, and hence $P_3^{n+1,n+1}$
is not Eulerian.  Similarly, $P_3^{2n-1,2n-1}(n-1)$ contains the two necklaces. $[2^{n-1},1]$ and
$[2^{n-2},0,0]$. These clearly do not share a vertex, and hence $P_3^{2n-1,2n-1}(n-1)$ is not
Eulerian.

Next, suppose $k$ is odd, $n+1<r<(k-1)n-1$ and $r\notin\{n-1+k/2,(k-1)(n-1)+k/2\}$.  In this case
$P_k^{r,r}(n-1)$ will always contain more than one necklace. However, no two necklaces can share a
vertex since $w_p$ is one-to-one when $k$ is odd, and hence a vertex can only have at most one
incoming or outgoing edge with any given pseudoweight.  Hence, $P_k^{r,r}(n-1)$ is not strongly
connected.

Finally, suppose $k$ is even and $n+1< r< (k-1)n-1$.  Suppose $r=n+s(k-2)+t-1$, where $1\leq t<k-1$
and $0\leq s<n$.  Then there exists an edge in $P_k^{r,r}(n-1)$ of the form
$((k-1)^s,t,1^{n-s-1})$. No vertex in this necklace will have more than one incoming edge or
outgoing edge in $P_k^{r,r}(n-1)$ except in the case $t=k/2$.  If $t=k/2$ then the necklace
$[(k-1)^s,0,1^{n-s-1}]$ will share a vertex with the necklace $[(k-1)^s,t,1^{n-s-1}]$ but neither
necklace will share a vertex with any other edges in $P_k^{r,r}(n-1)$.

We next show that $P_k^{r,r}(n-1)$ will always contain another necklace not equal to either
$[(k-1)^s,t,1^{n-s-1}]$ (or $[(k-1)^s,0,1^{n-s-1}]$ if $t=k/2$), establishing that $P_k^{r,r}(n-1)$
is not strongly connected. We need to consider three cases.
\begin{itemize}
\item First suppose $s=0$, i.e.\ there exists an edge in $P_k^{r,r}(n-1)$ of the form
    $(t,1^{n-1})$, where $t>2$ since $r>n+1$.  Then there will also exist an edge of the form
    $(t-1,2,1^{n-2})$ in $P_k^{r,r}(n-1)$, since $n\geq2$, i.e.\ another necklace.
\item Next suppose $s>0$ and $t<k-2$; then, since $k\geq4$, there exists a further edge of the
    form $((k-1)^{s-1},k-2,t+1,1^{n-s-1})$ in $P_k^{r,r}(n-1)$.
\item Finally suppose $s>0$ and $t=k-2$.  Now, since $r< (k-1)n-1$, $n-s-1>0$.  Then the edge
    mentioned above takes the form $((k-1)^s,k-2,1^{n-s-1})$.  An additional edge is then
    $((k-1)^s,k-3,2,1^{n-s-2})$.
\end{itemize}
The result follows.
\end{proof}

\end{document}
